\chapter{Le SOC : Wazuh sur le lab}
\label{ch:soc-wazuh}

Ce chapitre décrit le SOC du lab~: comment Wazuh est déployé, ce
qu'il surveille, et ce qu'il a détecté pendant la période de
tests. C'est aussi la principale source de preuves pour le
dossier d'homologation~: les règles actives, les agents
connectés, les alertes générées sont autant d'éléments que
HOMO-CI va collecter.

\section{Pourquoi Wazuh}

Trois raisons~:

\begin{enumerate}
  \item \textbf{C'est open source}, donc auditable, modifiable,
    sans coût de licence. Compatible avec le crédit étudiant
    Azure.
  \item \textbf{C'est tout-en-un.} Manager, indexer (basé sur
    OpenSearch), dashboard~: tout est dans une seule VM en
    déploiement all-in-one. Pas besoin d'orchestrer plusieurs
    machines pour un lab.
  \item \textbf{La base de règles est très fournie}~: par défaut,
    Wazuh livre plusieurs centaines de règles couvrant SSH,
    Windows Event Log, Docker, nginx, Apache, PostgreSQL,
    Active Directory.
\end{enumerate}

\section{Installation et version}

Wazuh est installé en mode \emph{all-in-one} sur \texttt{vm-siem01}
via le script officiel~:

\begin{lstlisting}[language=bash,caption={Installation Wazuh
all-in-one (extrait).}]
curl -sO https://packages.wazuh.com/4.9/wazuh-install.sh
sudo bash wazuh-install.sh -a
\end{lstlisting}

Version installée~: \textbf{Wazuh 4.9.2-1}. Le pinning de la
version est important~: le repository \texttt{4.x} sert toujours
la dernière minor, et le manager 4.9.2 refuse les agents 4.10.x
(message d'erreur peu clair lors de l'enrôlement). Le pin se fait
côté agent au moment du \texttt{apt install}~:

\begin{lstlisting}[language=bash,caption={Installation d'un agent
sur Ubuntu.}]
sudo apt install wazuh-agent=4.9.2-1
\end{lstlisting}

\section{Agents déployés}

\begin{table}[h]
\centering
\caption{Agents Wazuh dans le lab.}
\label{tab:agents}
\small
\begin{tabular}{lll}
\toprule
\textbf{Hôte} & \textbf{OS} & \textbf{Sources surveillées} \\
\midrule
vm-web01 & Ubuntu 22.04 & syslog, auth.log, nginx, Docker, apps PHP \\
vm-dc01 & Windows Server 2022 & Security event log, Sysmon, AD-related \\
\bottomrule
\end{tabular}
\end{table}

Le SIEM lui-même (\texttt{vm-siem01}) n'a pas d'agent~: il s'auto-
surveille via ses propres logs internes.

\section{Sources de logs côté serveur web}

L'agent sur \texttt{vm-web01} envoie à Wazuh~:

\begin{itemize}
  \item \texttt{/var/log/auth.log} et \texttt{/var/log/syslog}
    (SSH, sudo, cron, systemd).
  \item \texttt{/var/log/nginx/access.log} et
    \texttt{/var/log/nginx/error.log} (trafic HTTP/HTTPS).
  \item \texttt{/var/log/corpnet/internal/maformation/maformation.log}
    et \texttt{macandidature/macandidature.log} (logs applicatifs
    JSON).
  \item Les événements Docker via le module
    \texttt{docker-listener} (création/arrêt/exec dans les
    containers).
\end{itemize}

\section{Sources de logs côté AD}

Sur \texttt{vm-dc01}, l'agent Wazuh récupère le
\emph{Security Event Log} Windows. Les événements clés
surveillés~:

\begin{itemize}
  \item \textbf{4624} (logon) / \textbf{4625} (échec de logon)~:
    détection de brute force.
  \item \textbf{4768} / \textbf{4769} / \textbf{4771} (Kerberos)~:
    détection de \emph{Kerberoasting}, AS-REP Roasting.
  \item \textbf{4720} / \textbf{4732}~: création d'utilisateur,
    ajout à un groupe privilégié.
  \item \textbf{4662}~: opérations sur les objets AD (utile pour
    détecter une tentative de \emph{DCSync}).
\end{itemize}

\section{Règles de détection : built-in et custom}

Wazuh fournit plusieurs centaines de règles par défaut. Celles qui
n'étaient pas couvertes ont été ajoutées sous forme de fichiers
\texttt{local\_rules.xml}~:

\begin{itemize}
  \item Détection des injections SQL dans les logs nginx (motifs
    \texttt{UNION SELECT}, \texttt{OR 1=1}, etc.).
  \item Détection XSS sur les paramètres POST (balises
    \texttt{<script>}, événements \texttt{onerror}).
  \item Détection IDOR via énumération de paramètres
    \texttt{conversation\_id}.
  \item Détection \texttt{docker exec} initié depuis un container
    web (corrélation entre l'événement Docker et le PID
    parent).
  \item Détection \emph{Kerberoasting} (événement 4769 avec
    chiffrement RC4 sur un compte de service).
\end{itemize}

\section{Ce que le SOC a vu pendant les tests}

Sur la période d'exercices (scénarios d'attaque rejoués pour
valider la détection), Wazuh a généré \textbf{927 alertes}.
Décomposition~:

\begin{table}[h]
\centering
\caption{Alertes Wazuh par catégorie (période de tests).}
\label{tab:alertes}
\small
\begin{tabular}{lr}
\toprule
\textbf{Catégorie} & \textbf{Nombre} \\
\midrule
SSH brute force / scans & 412 \\
Attaques web (XSS, SQLi, IDOR) & 287 \\
AD (Kerberoasting, énumération) & 156 \\
Docker (exec suspect, escape) & 38 \\
Autres (sudo, file integrity) & 34 \\
\midrule
\textbf{Total} & \textbf{927} \\
\bottomrule
\end{tabular}
\end{table}

Ces chiffres ne sont pas un résultat scientifique~: ils dépendent
des scénarios joués et de leur fréquence. Mais ils valident que
la chaîne (agent~$\to$~manager~$\to$~indexer~$\to$~dashboard) fonctionne et
que les règles déclenchent comme attendu.

\section{Active Response : du bruit à la réaction}
\label{sec:active-response}

Une fois les détections en place, l'étape suivante était de
mesurer si Wazuh pouvait \emph{réagir} et pas seulement
\emph{alerter}. On a activé le mécanisme \emph{Active Response}
sur un scope contrôlé~: déclenchement uniquement sur les règles
de brute force (SSH et applicatif), action \texttt{firewall-drop}
(ajout d'une règle \texttt{iptables DROP} sur l'agent attaqué),
durée du blocage 600~s puis débanissement automatique.

\paragraph{Configuration.} Deux blocs \texttt{<active-response>}
ajoutés à \texttt{/var/ossec/etc/ossec.conf} sur
\texttt{vm-siem01}~:

\begin{lstlisting}[language=XML,caption={Réponse active sur
brute force SSH et applicatif.}]
<active-response>
  <command>firewall-drop</command>
  <location>local</location>
  <rules_id>5712</rules_id>   <!-- SSH brute force built-in -->
  <timeout>600</timeout>
</active-response>

<active-response>
  <command>firewall-drop</command>
  <location>local</location>
  <rules_id>100302</rules_id>  <!-- login brute force apps -->
  <timeout>600</timeout>
</active-response>
\end{lstlisting}

\paragraph{Scénario de test.} Depuis le manager (qui partage le
réseau privé Azure avec \texttt{vm-web01}), 12 tentatives SSH
avec un utilisateur inexistant (\texttt{rootbf}) et un mot de
passe erroné, lancées en rafale. La règle 5712 (built-in Wazuh)
déclenche à 8~échecs en moins de 120~s, niveau~10.

\paragraph{Chronométrage observé.}

\begin{table}[h]
\centering
\caption{Latence détection $\to$ blocage SSH brute force (mesure
unique, 16~mai~2026).}
\label{tab:active-response}
\small
\begin{tabular}{lrr}
\toprule
\textbf{Étape} & \textbf{Horodatage UTC} & \textbf{$\Delta$ depuis $t_0$} \\
\midrule
1\textsuperscript{re} tentative SSH ($t_0$) & 09:47:47.799 & 0\,s \\
8\textsuperscript{e} tentative (seuil 5712) & 09:47:48.360 & 0,56\,s \\
Règle 5712 déclenchée au manager & 09:47:50.254 & \textbf{2,45\,s} \\
\texttt{firewall-drop} exécuté sur web01 & 09:47:50 & $\sim$2,5\,s \\
Sonde TCP en \emph{timeout} (blocage confirmé) & 09:47:55+ & bloqué \\
\bottomrule
\end{tabular}
\end{table}

Vérification côté \texttt{vm-web01}~:

\begin{lstlisting}[language=bash,caption={État iptables après
déclenchement.}]
$ sudo iptables -L INPUT -n --line-numbers
Chain INPUT (policy ACCEPT)
num  target  prot  source       destination
1    DROP    all   10.0.3.10    0.0.0.0/0
\end{lstlisting}

La règle disparaît automatiquement après 600~s. Le journal
\texttt{/var/ossec/logs/active-responses.log} sur l'agent
contient le payload JSON complet de l'alerte qui a déclenché
l'action (rule~id, MITRE T1110, srcip, timestamp), ce qui
fournit une preuve auditable pour HOMO-CI.

\paragraph{Périmètre volontairement restreint.} L'Active Response
est limitée à deux règles de brute force. On n'a pas activé
\texttt{disable-account} (qui désactive le compte ciblé) ni
\texttt{host-deny} (entrée \texttt{/etc/hosts.deny}), pour
éviter qu'une fausse alerte ne bloque un opérateur légitime
pendant les tests. C'est un choix conservateur~: en production
sur un vrai ENT, on déciderait du périmètre en concertation avec
le RSSI.

\section{Ce qu'il reste à faire côté SOC}

Pour rester honnête sur le périmètre~:

\begin{itemize}
  \item \textbf{Pas de tuning avancé des règles}~: il y a
    probablement des faux positifs sur les 927~alertes mesurées
    qu'on n'a pas filtrés. Un vrai déploiement passerait par une
    phase de \emph{baselining} et l'écriture d'exceptions ciblées
    (rule \texttt{<if\_sid>} avec \texttt{<options>no\_full\_log</options>}
    pour les sources connues).
  \item \textbf{Pas d'enrichissement avec des IOC externes}~:
    MISP, OpenCTI, ou des feeds publics (URLhaus, OTX) seraient
    branchés via le module \texttt{integrator} de Wazuh. C'est
    une perspective documentée dans le chapitre~\ref{ch:limites},
    pas un blocage technique.
\end{itemize}

\begin{takeaway}
Le SOC du lab est volontairement simple~: Wazuh 4.9.2 all-in-one,
deux agents, règles built-in + une dizaine de règles custom,
Active Response activée sur les brute force. C'est suffisant
pour générer des preuves d'audit exploitables par HOMO-CI (état
des agents, règles actives, alertes par catégorie, et désormais
événements d'\emph{Active Response} datés).
\end{takeaway}
